\clearpage
\fancyhead[L]{\small ZACH. L'ATTRACTION DES MONTAGNES.}
\hypertarget{native-574-lower}{}
\begin{center}\small Göttingische gelehrte Anzeigen. No. 210. Pages 2089..2100. 31 December 1814.\\[3pt]\rule{.72\linewidth}{.3pt}\end{center}

\emph{L'attraction des montagnes, et ses effets sur les fils à plomb ou sur les niveaux des
instrumens d'astronomie, constatés et determinés par des observations astronomiques et
géodésiques faites en} 1810, \emph{à l'ermitage de notre dame des anges, sur le mont de Mimet, et
au fanal de l'isle de Planier près de Marseille; suivis de la description géométrique de la
ville de Marseille et de son territoire. Par le baron} \textsc{de Zach}. 1814.
Avignon. Published by \textsc{Seguin} the elder. XX and 716 pages in large octavo.

The observations made by \textsc{Bouguer} and \textsc{Condamine} in 1738 on Chimborasso and by \textsc{Maskelyne} in 1774 on Shehallien in Scotland, on the effect of large mountain masses on the direction of gravity, had so far been the only direct experiments on this interesting subject. The former gave $7''5$ for the deflection of the plumb line from its natural direction, the latter $5''8$. Strictly speaking, however, only the observations in Scotland can be considered decidedly reliable. The observations in Peru certainly cannot, since the conditions were too unfavourable and the instruments comparatively too imperfect. Given the importance of the subject and the greatly improved state of practical astronomy, it is highly desirable that more such interesting experiments be made. The present work by one of the foremost observers must therefore be regarded as a genuine gain for science. The work under review contains every observation belonging to this undertaking down to the smallest detail, sets out the methods of observation and calculation at length, and also gives numerous carefully calculated auxiliary tables that can be used to great advantage elsewhere. In addition it offers a wealth of practical remarks and judgments of many kinds, as well as scholarly notices and subsidiary investigations, a merit that readers have come to expect in the author’s writings.

\SourceLeaf{575}In the introduction the author first leads the reader to the point of view from which the phenomenon of mountain attraction and the observations that reveal it must be considered. He then describes the relevant experiments made so far and finally gives a general account of his own operations. It may suffice here to recall that gravity is simply the combined effect that all the constituent particles of the Earth exert, under the general law of attraction, on bodies at its surface, as modified by the centrifugal force arising from the Earth’s rotation. Astronomical latitude, as the angle that the direction of gravity makes with the plane of the equator, must therefore increase regularly from the equator toward the pole, following laws derived by comparing the various measurements of meridian arcs, insofar as the Earth is regarded as a regular body. This regular progression is disturbed, however, where large departures of the Earth’s surface from its normal shape cause perceptible deflections of gravity from its natural direction. Consider two places. At one of them a nearby mountain mass affects the direction of gravity from north to south or from south to north, while at the other the influence of that mass either is no longer perceptible or acts in the opposite direction. The difference of their latitudes as found by astronomical observations will differ from the result calculated from the relative position and distance of the two places. This discrepancy gives us the magnitude of the deflection in the first case or the sum of the two deflections in the other. Clearly, then, the second case, which arose in \textsc{Maskelyne}’s observations, is the more favourable. In the first the task is twice as difficult, since the whole deflection always amounts to only a very few seconds of arc and the observations must therefore be of the utmost precision. The terrain where Mr.\ \textsc{von Zach} carried out his operations probably did not permit him to use that more favourable method. One of his observing stations was on the southern slope of Mount Mimet, the highest point (400 toises above sea level) in a range of limestone mountains running from east to west, about two Meilen north of Marseille. The other was the lighthouse on the very small island of Planier, a couple of Meilen southwest of Marseille. At the first site, at an elevation of 250 toises near the ruined monastery of Notre-Dame des Anges, the mountain had to deflect the direction of gravity northward and thus reduce the latitude. At the other site, because of the great distance, the effect of the mountain’s attraction could be regarded as imperceptible. The observed astronomical latitudes therefore had to give a smaller difference than the calculation from the positions of the two places as determined by geodetic measurements, on the assumption of a regular progression. In determining the latitudes it is customary to use the same stars at both places. In this way the difference in latitude, which alone matters here, becomes independent of the absolute accuracy of the stars’ declinations. One can even leave the absolute latitudes out of account and use, in place of their difference, the difference between the observed zenith distances of the stars after removing the small periodic inequalities and reducing the distances to the same instant.

Of the eight sections into which Mr.\ \textsc{von Zach} has divided the work, the first two contain all the astronomical observations made at Notre-Dame des Anges and on the island of Planier. The observations were of three kinds. There were observations of the zenith distances of the three stars $\alpha$ Ophiuchus, $\zeta$ Aquila and $\alpha$ Aquila, of gunpowder signals to determine the differences in longitude relative to the Marseille observatory, and of azimuths to orient the network of triangles linking the two observing sites. The times \SourceLeaf{576}were determined with three \textsc{Emery} chronometers using corresponding altitudes of the Sun. The zenith distances were observed with a twelve-inch repeating circle by \textsc{Reichenbach}, and the azimuths with an eight-inch theodolite by the same instrument maker. The general reduction table given here for zenith distances observed away from culmination is valuable to observers using repeating circles, since it has been calculated more precisely and more completely than those available elsewhere. The third section contains the terrestrial measurements. A baseline of 1182,4 toises was measured along the road from Marseille to Aix with wooden measuring rods, whose small variability the author illustrates here with noteworthy observations. The horizontal angles of the seven triangles in the network connecting the two main points to the baseline were also observed with the theodolite. Most were repeated ten times and some more often. The largest deviation in the sum of the three angles of a triangle was $4''9$. In the fourth section the author derives the differences in latitude and longitude between the two observing sites from these geodetic measurements. He takes as his basis the spheroidal shape of the Earth, the flattening $\frac{1}{310}$ and an equatorial radius of 3271604 toises, and uses the formulas given by \textsc{Delambre} in Méthodes analytiques pour la détermination d'un arc du méridien. He finds $12'\,3''11$ for the former difference and $15'\,46''46$ for the latter. The author also provides several auxiliary tables for these calculations. In the fifth section these calculated results are compared with those of the astronomical observations. The observations gave the following differences in latitude.

\begin{center}\small\begin{tabular}{@{}l@{\hspace{9pt}}r@{}}
598 observations of $\alpha$ Ophiuchus $.\;.\;.\;.$ & $12'\,0''84$ \\
518 observations of $\zeta$ Aquila $.\;.\;.\;.\;.$ & $12\phantom{'}\,1.26$ \\
654 observations of $\alpha$ Aquila $.\;.\;.\;.\;.$ & $12\phantom{'}\,1.30$ \\
\end{tabular}\end{center}
\noindent Thus the mean of 1770 observations gave $12'\,1''13$, so that $1''98$ is to be regarded as the final result for the effect of the attraction of Mimet at the observing site of Notre-Dame des Anges. Observations of the gunpowder signals had yielded the following differences in longitude between the Marseille observatory and the two sites.

\begin{center}\small\begin{tabular}{@{}l@{\hspace{9pt}}r@{}}
Notre-Dame des Anges $.\;.\;.\;.\;.\;.\;.\;.$ & $7'\,29''25$ \\
Lighthouse on the island of Planier $.\;.\;.\;.$ & $8\phantom{'}\,5.70$ \\
\end{tabular}\end{center}
\noindent The former came from 63 observations on 11 different days, the latter from 53 observations on 12 days (with five observations from one day excluded solely because they gave a result, $7'\,38''1$, that differed too much from the others). The total difference in longitude between Notre-Dame des Anges and the island of Planier is therefore $15'\,34''95$ from these observations, or $15'\,35''79$ when reduced to the points to which the geodetic measurements refer. It is thus $10''67$ smaller than the result of the geodetic measurements, or $13''05$ smaller if no observations at all are excluded. There is no justification for attributing this difference to the mountain’s attraction as well. It must instead be attributed, at least for the most part, to the unavoidable observational errors in the astronomical determination of the difference in longitude. That determination depended on time and could come nowhere near the accuracy of the observations of zenith distances. In any case, the entire longitude determination was a secondary operation, quite inessential to the main undertaking. It is enough that it shows us what errors even the most experienced observer is still liable to encounter with the method used. Comparisons of the different observed azimuths offer a similar lesson. When the azimuths determined on the island of Planier are transferred to Notre-Dame des Anges, the differences between \SourceLeaf{577}the azimuths thus derived and those observed directly there are 7 and 21 seconds of arc. These discrepancies can readily be explained by apportioning them fairly between the azimuth determinations and the terrestrial angle observations. We cannot refrain, however, from mentioning a small detail that matters precisely in comparisons of \emph{this} kind. We noticed that in centring the angles, the distances of the observing sites from the triangulation points, which the author gives in ten-thousandths of a toise, come out, when converted to metric measure, almost all to exact tenths of a metre, that is, whole decimetres. This hardly seems a coincidence. The author appears to have measured those distances with a metric rule and disregarded smaller divisions. Yet for several angles an error of one decimetre (about 4 inches) changes the reduction by $3''$. Admittedly, this is of very little consequence for the author’s main purpose. But if the transfer of an azimuth to another, distant place and its comparison with an azimuth observed directly there are to yield new information about the shape of the Earth, the greater the number of intermediate triangles, the more care one must take in this process of centring. The author himself warmly recommends carrying out such an operation with \textsc{Reichenbach} instruments and proposes the Lipari Islands for it. We wish for nothing more than to see the author carry out this undertaking himself, which we have all the more reason to hope for now that he has changed his place of residence.

The sixth section is the most substantial in the whole work. The author first examines how errors that might occur in each operation could affect the final result. It is easy to see that observing the zenith distances is by far the most delicate part of the whole undertaking. Given how exceptionally small the final result is, and that its reliability requires those observations to be accurate even to fractions of a second of arc, some might have all the more reason for reservations because the author himself said on an earlier occasion that every absolute latitude determined by a repeating circle remains subject to an uncertainty of several seconds. That statement was admittedly based on experience showing such discrepancies between determinations made with \emph{different} circles, whereas the results obtained with any one circle agreed admirably among themselves. In the present undertaking, by contrast, all observations were made with one and the same instrument, they agree with each other extremely well, and only small differences matter here, independently of the absolute determinations themselves. To corroborate his results further, Mr.\ \textsc{von Zach} uses the numerous astronomical observations he made partly at the Royal Observatory in Marseille and partly at his own observatories, La Capellette and St.\ Peyre. The differences between the latitudes thus determined for those three points (which are no longer appreciably affected by the attraction of the mountains) and the latitude of the island of Planier agree extremely closely with the results of the geodetic measurements by which he connected them to his network of triangles. A similar comparison with Notre-Dame des Anges, however, again gives that difference of two seconds with fairly consistent results. One might perhaps question whether this is admissible as corroboration, since the five latitudes are based on different stars (Notre-Dame des Anges and the island of Planier on $\alpha$ Aquila, and the other three on $\alpha$ and $\beta$ in the Little Bear). But if one gives due weight to the fact that the declinations used are precisely those that Mr.\ \textsc{von Zach} himself determined with this very instrument at Milan, a place not very much farther north, one readily sees that the doubt is groundless, provided one accepts the assumption that \SourceLeaf{578}any errors there may be in the absolute zenith distances determined with the circle are constant for any given zenith distance and nearly equal for zenith distances that differ little. Whatever one may think of the causes of the remarkable discrepancies that Mr.\ \textsc{von Zach} found with \emph{different} circles, the persistent agreement of the results obtained with the same circle cannot, it must be admitted, be explained without that assumption. For all these reasons one may take the difference of 2 seconds that was found to represent the attraction of Mount Mimet very closely indeed. Yet no practical astronomer will deny that the result still has an uncertainty of half a second (if not more), so that it cannot be confined within narrower limits than $1''5$ to $2''5$. One must therefore regret that so fine and so carefully executed an undertaking yielded so small a quantity, with an unavoidable uncertainty so large in proportion to it. It is all the more regrettable that the second site was chosen on the island of Planier instead of on the north side, right next to Mount Mimet. Mr.\ \textsc{von Zach} does not say what kept him from doing so. But so discerning an astronomer would surely not have forgone an advantage that might have doubled the effect, had the terrain not presented insuperable obstacles. If the effect of Mimet had turned out to be larger, we would also have very much wished for a fuller account of the physical dimensions of the mountain. Under the circumstances, however, any conclusions drawn from them about the relative densities of Minet and the whole Earth would be of too limited an accuracy. Even so, Mr.\ \textsc{von Zach}’s measurements have the important benefit of proving at least how extraordinarily small the influence of so substantial a mountain mass is, making us a little more wary of invoking the possible influence of local attractions too hastily when measurements refuse to agree.

Besides the investigations mentioned, the sixth section examines some determinations by \textsc{Cassini de Thury} in Méridienne vérifiée which differ appreciably from those of the author. Far greater errors, amounting in fact to distortions, are found in \textsc{Cassini}’s description géometrique de la France and in his large map of France. --- Finally, the author also surveys the results for the positions of several stars, drawn partly from his present work, partly from his own earlier observations and partly from a critical discussion of observations by other astronomers. Here we mention only the right ascension of the pole star, which is so important in many respects and for which \textsc{Mathieu}’s observations in 1812 gave an increase of 5 seconds of time. To bring \textsc{von Zach}’s determination of 1790 into agreement with this, one must assume a proper motion. Mr.\ \textsc{von Zach} finds it fully confirmed by comparison with \textsc{Lacaille}’s observations of 1750 and puts it at $+3''177$ of arc annually. (The most precise determination of the current right ascension of the pole star will come from the numerous observations Mr.\ \textsc{von Lindenau} has made at the Seeberg observatory over the past five years. We can already say, provisionally, that they confirm the increase found by \textsc{Mathieu}.)

The limited space in our journal allows us only to touch briefly on the remaining sections of \textsc{von Zach}’s work. The seventh section determines the heights of the triangulation points and of some other places above the Mediterranean Sea. The eighth section gives the positions of a great many points in the city of Marseille and its surroundings by their distances from the meridian and perpendicular of the Royal Observatory there and by their longitude \SourceLeaf{579}and latitude. It also contains several other interesting critical investigations of various sites in Marseille that have become noteworthy in the history of astronomy. Finally, in an appendix, Mr.\ \textsc{von Zach} gives a new calculation of all the astronomical observations that \textsc{Maskelyne} made at Shehallien, which leaves the final result found by \textsc{Maskelyne} himself unchanged.
